\documentclass[../../../include/open-logic-section]{subfiles}

\begin{document}

\olfileid{sth}{choice}{intro}

\olsection{ভূমিকা}

\crefrange{sth:cardinals::chap}{sth:card-arithmetic::chap}-এ
অঙ্কবাচক সংখ্যাকে ক্রমসংখ্যা ধরে তার তত্ত্ব গড়েছি।
এই পদ্ধতি সুক্রমবিন্যাসের স্বতঃসিদ্ধের উপর নির্ভর করে।
দেখা যাবে, সুক্রমবিন্যাস আরেকটি নীতির—চয়নের
স্বতঃসিদ্ধের—সমতুল্য। এই স্বতঃসিদ্ধ গ্রহণযোগ্য কি না,
তা নিয়ে গুরুতর দার্শনিক আলোচনা হয়েছে। এই অধ্যায়ের
প্রশ্ন দুটি: স্বতঃসিদ্ধটি কীভাবে কাজে লাগে, এবং তার
ন্যায্যতা কি দেখানো যায়?

\end{document}
